\section{A regular coefficient for the inverse terminal operator}
\label{rl-section}

We next reconstruct an actual full-group coefficient used in the
four-operator Fourier argument.  Radial coordinates involve ratios
of consecutive minors, which may be undefined inside the relevant
open subset of the full group.  The construction below has only the
positive Gram polynomial in its denominator.  It gives the
zero-order spectral factor used in the source
\cite{alpo-uniform-source}, including covariance, growth of every
derivative, and preservation of the coefficient class.

\subsection{The Hermitian space, minors and cofactor vector}

Let $k\geq1$ and $r\geq0$.  Fix a Witt decomposition
\[
 V=U\oplus W\oplus U^\vee,\qquad
 U=\langle e_1,\ldots,e_r\rangle,\quad
 U^\vee=\langle e_1^\vee,\ldots,e_r^\vee\rangle,
\]
for a Hermitian form $h$, linear in its first argument.
The subspaces $U,U^\vee$ are isotropic,
$h(e_i,e_j^\vee)=\delta_{ij}$, and $W$ is perpendicular
to their sum and has signature $(k,1)$.
Choose a basis of $W$ in which the form is
$\operatorname{diag}(1,\ldots,1,-1)$.
In the resulting basis of $V$ its matrix $H$ satisfies
$H^*=H$ and $H^2=1$.
Write $\mathbf G=U(V,h)$.

Fix $z,z'\in W$ with
\[
                 h(z,z)=-T,\qquad h(z',z')=-T',
                         \qquad T,T'>0.
\]
Put $f_i=e_i$, $f'_i=e_i$ for $1\leq i\leq r$, and
$f_{r+1}=z$, $f'_{r+1}=z'$.
For $g\in\mathbf G$ define
\begin{align}
 Z(g)_{ab}&=h(gf'_b,f_a),\qquad 1\leq a,b\leq r+1,
                                      \label{rl-Z}\\
 D(g)&=\det Z(g),\qquad
 d(g)=\det Z(g)_{\{1,\ldots,r\},\{1,\ldots,r\}},\nonumber\\
 \mathcal G(g)&=|D(g)|^2-TT'|d(g)|^2. \label{rl-Gram}
\end{align}
When $r=0$, $d(g)=1$.
These are the last two minors and the actual Gram polynomial
in the source, before passage to radial coordinates.
Define
\begin{equation}\label{rl-cofactor-vector}
 L(g)=\sum_{a=1}^{r+1}
        \overline{\operatorname{adj}(Z(g))_{r+1,a}}\,f_a,
 \qquad y(g)=g^{-1}L(g).
\end{equation}
Let $\omega(g)$ be the projection of the $W$-component
of $y(g)$ onto the positive space
$W'_+=W\cap(z')^\perp$.
Fix an orthonormal basis $B_1,\ldots,B_k$ of $W'_+$ and put
\[
                 \Omega_j(g)=h(y(g),B_j).
\]
Thus $\omega(g)=\sum_j\Omega_j(g)B_j$.

\begin{lemma}\label{rl-cofactor-norm}
For every $g\in\mathbf G$,
\begin{align}
 h(y(g),f'_b)&=\overline{D(g)}\,\delta_{b,r+1},
                                                  \label{rl-dual-pairing}\\
 h(y(g),y(g))&=-T|d(g)|^2,\label{rl-y-norm}\\
 \sum_{j=1}^k|\Omega_j(g)|^2
             &=\frac{\mathcal G(g)}{T'}.\label{rl-omega-norm}
\end{align}
In particular $\mathcal G(g)\geq0$.
Each $\Omega_j$ is the restriction to $\mathbf G$ of a
polynomial of degree $r+1$ in the conjugates of the
matrix entries of $g$.
No minor is required to be nonzero in these assertions.
\end{lemma}

\begin{proof}
Since $g$ is unitary,
\[
 h(y,f'_b)=h(L,gf'_b)
   =\sum_a
     \overline{\operatorname{adj}(Z)_{r+1,a}}\,
     \overline{Z_{ab}}
   =\overline{D}\,\delta_{b,r+1}.
\]
This is the adjugate identity, valid also when $Z$ is singular.
All $f_a$ with $a\leq r$ are mutually perpendicular
and perpendicular to $z$, so
\[
 h(y,y)=h(L,L)
       =-T|\operatorname{adj}(Z)_{r+1,r+1}|^2
       =-T|d|^2.
\]
The first $r$ equations in \eqref{rl-dual-pairing}
say that $y\in U^\perp=U\oplus W$.
Its component along $z'$ is
$-\overline D\,z'/T'$.  Consequently
\[
           y=u-\frac{\overline D}{T'}z'+\omega
                          \qquad(u\in U,\ \omega\in W'_+).
\]
Taking its norm gives
$-T|d|^2=-|D|^2/T'+h(\omega,\omega)$,
which proves \eqref{rl-omega-norm}.

For the polynomial assertion, $Z$ is linear in the
entries of $g$, so every cofactor has degree $r$.
On $\mathbf G$ one has $g^{-1}=Hg^*H$.
Replacing $g^{-1}$ by this formula in
\eqref{rl-cofactor-vector} gives a polynomial of
degree $r+1$ in the conjugate entries.
Pairing with the fixed $B_j$ is linear and does not
change that degree.
\end{proof}

\subsection{Covariance, including the unipotent subgroups}

Let $N\subset\mathbf G$ consist of the elements satisfying
\[
 n e_i-e_i\in\langle e_1,\ldots,e_{i-1}\rangle,
       \qquad n w-w\in U\quad(w\in W).
\]
This is the source's upper unipotent subgroup in the
chosen Witt basis.  It preserves $U^\perp$ and acts
as the identity on $U^\perp/U$.  The same assertions
hold for its inverse, by triangular inversion.
Let $K_z$ and $K_{z'}$ fix every $e_i,e_i^\vee$
and fix $z$ and $z'$, respectively.  They act on
$W$ and are the compact groups isomorphic to $U(k)$.

\begin{lemma}\label{rl-cofactor-covariance}
Let $h_L\in K_z$, $h_R\in K_{z'}$ and
$n_L,n_R\in N$, and set
\[
                  g'=h_L^{-1}n_L^{-1}g n_Rh_R.
\]
Every leading principal minor of $Z(g')$ equals the
corresponding minor of $Z(g)$.  Moreover
\begin{equation}\label{rl-omega-covariance}
                \omega(g')=h_R^{-1}\omega(g),
       \qquad \mathcal G(g')=\mathcal G(g).
\end{equation}
These statements hold also when intermediate minors vanish.
\end{lemma}

\begin{proof}
Write $f$ for the matrix with columns $f_a$, and
similarly $f'$ for $f'_b$.  There are upper triangular
unipotent matrices $A,B$ of size $r+1$ with
$n_Lf=fA$ and $n_Rf'=f'B$.
Since the compact factors fix these columns,
unitarity gives
\[
                         Z(g')=A^*Z(g)B.
\]
For a leading block, upper triangularity of $A,B$
means that multiplication uses only the corresponding
leading blocks.  Their determinants are one.
This proves the minor identities, and hence the
identity for $\mathcal G$.

For completeness, the cofactor transformation can be
obtained without assuming that $Z$ is invertible.
The identity
\[
 \operatorname{adj}(A^* Z B)
       =B^{-1}\operatorname{adj}(Z)(A^*)^{-1}
\]
holds for invertible $Z$ by the inverse formula;
both sides are polynomial in the entries of $Z$,
so it holds for every $Z$ by continuity and the
density of invertible matrices.  The last row of
$B^{-1}$ is the last coordinate row.  Transposing
and conjugating that row identity in
\eqref{rl-cofactor-vector} yields
\[
                 L(g')=n_L^{-1}L(g)
                       =h_L^{-1}n_L^{-1}L(g).
\]
The second equality holds because
$n_L^{-1}L(g)$ lies in $\langle f_1,\ldots,f_{r+1}\rangle$,
which $h_L$ fixes pointwise.  Thus
$y(g')=h_R^{-1}n_R^{-1}y(g)$.
The vector $y(g)$ belongs to $U^\perp$ by
Lemma~\ref{rl-cofactor-norm}, and $n_R^{-1}$
does not change its $W$-component.
Finally $h_R$ preserves $W$, fixes $z'$, and
acts unitarily on $W'_+$.  Projecting onto $W'_+$
proves \eqref{rl-omega-covariance}.
\end{proof}

\subsection{The radial cell and the compact Capelli contraction}

For $c=(c_1,\ldots,c_r)\in(\mathbb C^\times)^r$,
let $b(c)\in\mathbf G$ fix $W$ and satisfy
\[
 b(c)e_i=c_i e_i^\vee,\qquad
 b(c)e_i^\vee=\overline c_i^{-1}e_i.
\]
Let $m\in U(W,h|_W)$, extended as the identity on
$U\oplus U^\vee$.  Put
\[
             a(c)=\prod_{i=1}^r c_i,\qquad
             \zeta(m)=h(mz',z),\qquad
             G(m)=|\zeta(m)|^2-TT'.
\]

\begin{lemma}\label{rl-cell}
For $g=b(c)m$ one has
\begin{align}
 Z(g)&=\operatorname{diag}(c_1,\ldots,c_r,\zeta(m)),\nonumber\\
 \mathcal G(g)&=|a(c)|^2G(m),\label{rl-cell-Gram}\\
 \Omega_j(g)&=\overline{a(c)}\,h(m^{-1}z,B_j).
                                                 \label{rl-cell-omega}
\end{align}
\end{lemma}

\begin{proof}
Pair the images of the indicated basis vectors.
The decomposition into $W$ and its perpendicular
hyperbolic planes gives the diagonal matrix.
Its last cofactor row has only the last entry
$a(c)$, so $L(g)=\overline{a(c)}z$.
Since $b(c)$ fixes $z$, it follows that
$y(g)=\overline{a(c)}m^{-1}z$.
The remaining statements follow from the definitions.
\end{proof}

Let $\pi':K_{z'}\to U(W_{\pi'})$ be a finite-dimensional
continuous unitary representation.  Its complexified
derivative in the fixed basis $B_1,\ldots,B_k$
gives the operators $E_{ij}$ of
Lemma~\ref{rp-capelli}.  For a unit vector
$\nu=(\nu_1,\ldots,\nu_k)\in\mathbb C^k$ define
\begin{equation}\label{rl-cap-vector}
 C_\nu(Y)=\sum_{i,j=1}^k(-1)^{i+j}
       \overline{\nu_i}\nu_j\,
       Q_{\{1,\ldots,k\}\setminus\{i\},
         \{1,\ldots,k\}\setminus\{j\}}
          \left(Y+\frac{k-2}{2}\right).
\end{equation}
The indexing convention in the minor is that of
Lemma~\ref{rp-capelli}, with positions starting at zero
inside each ordered tuple.  The sign is unchanged
when the ambient indices are shifted from zero to one.

\begin{lemma}\label{rl-compressed-capelli}
The operator $C_\nu(Y)$ is the compact Capelli
determinant on the orthogonal complement of
$\sum_i\nu_iB_i$, with shift $Y+(k-2)/2$.
In particular it is independent of an orthonormal
basis of that complement, and
\[
 C_\nu(Y)^*=C_\nu(\overline Y).
\]
If $h\in K_{z'}$ and $\nu'$ is the coordinate vector
of $h^{-1}\sum_i\nu_iB_i$, then
\begin{equation}\label{rl-cap-covariance}
          C_{\nu'}(Y)=\pi'(h)^{-1}C_\nu(Y)\pi'(h).
\end{equation}
Also $C_{\lambda\nu}(Y)=C_\nu(Y)$ when $|\lambda|=1$.
\end{lemma}

\begin{proof}
Complete $\nu$ to a unitary matrix
$P=(v_1,\ldots,v_{k-1},\nu)$.
Apply all-minor covariance from Lemma~\ref{rp-capelli}
to the determinant on its first $k-1$ columns.
The complementary-minor formula gives
\[
 \det P_{\widehat i,\{1,\ldots,k-1\}}
                =(-1)^{i+k}\det P\,\overline{\nu_i}.
\]
The conjugate factor from the inverse matrix is
$(-1)^{j+k}\overline{\det P}\,\nu_j$.
Their product is
$(-1)^{i+j}\overline{\nu_i}\nu_j$, which proves
\eqref{rl-cap-vector} and its interpretation.
Frame independence and the adjoint formula follow
from the corresponding Capelli statements.
Conjugation by $h^{-1}$ carries the complement and
all its rank-one Lie algebra generators to the
new complement, giving \eqref{rl-cap-covariance}.
The scalar phase cancels directly in
\eqref{rl-cap-vector}.  When $k=1$ this is the
empty determinant $1$, so that case is included.
\end{proof}

On the open set $\mathbf G^+=\{g:\mathcal G(g)>0\}$,
define
\[
     \nu(g)=\sqrt{T'/\mathcal G(g)}\,\Omega(g),
                   \qquad C_g(Y)=C_{\nu(g)}(Y).
\]
Lemma~\ref{rl-cofactor-norm} makes $\nu(g)$ a unit vector.
In particular the square root is the positive real one.
Substitution into \eqref{rl-cap-vector} gives the
formula without square roots
\begin{equation}\label{rl-rational-matrix}
 C_g(Y)=\frac{T'}{\mathcal G(g)}
        \sum_{i,j=1}^k(-1)^{i+j}
            \overline{\Omega_i(g)}\Omega_j(g)
            Q_{\widehat i,\widehat j}
                   \left(Y+\frac{k-2}{2}\right).
\end{equation}
It is smooth on $\mathbf G^+$, including at the
zeros of intermediate principal minors.

\begin{proposition}\label{rl-actual-terminal-lift}
For a radial function $\Phi$ of Section~\ref{rd-section},
at $g=b(c)m$ with $G(m)>0$,
\begin{equation}\label{rl-terminal-match}
   G(m)^{-1}(\mathcal T_{22}(Y)\Phi)(c,m)
                              =\Phi(c,m)C_g(Y).
\end{equation}
Moreover, with $g'$ as in Lemma~\ref{rl-cofactor-covariance},
\[
                C_{g'}(Y)=\pi'(h_R)^{-1}C_g(Y)\pi'(h_R).
\]
\end{proposition}

\begin{proof}
By \eqref{rl-cell-omega}, $\nu(g)$ differs by the
unit scalar $\overline{a(c)}/|a(c)|$ from the
normalized projection of $m^{-1}z$ onto $W'_+$.
Its orthogonal complement in $W'_+$ is therefore
$\langle m^{-1}z,z'\rangle^\perp$.
Lemma~\ref{rl-compressed-capelli} and the actual terminal
identity of Corollary~\ref{rp-finite-terminal} prove
\eqref{rl-terminal-match}.
The covariance assertion follows from
\eqref{rl-omega-covariance} and
\eqref{rl-cap-covariance}.
\end{proof}

\subsection{The inverse polynomial and its regular full-group lift}

Retain the common finite type set $\mathcal S$ and
$p_\rho(Y)=\prod_{i=1}^{k-1}(Y+\rho_i)$ from
Section~\ref{rd-section}.  It is nonempty under the
interlacing hypotheses there.
Let $\sigma\in\{1,-1\}$; for the actual signed
source operator it is $\sigma=(-1)^{k-l_0}$.
Choose $Y\in\mathbb C$ such that $p_\rho(Y)\ne0$
for every $\rho\in\mathcal S$.  Define
\begin{align}
 P_Y(X)&=\prod_{\rho\in\mathcal S}
                         (X-\sigma p_\rho(Y)),\nonumber\\
 q_Y(X)&=\frac{1-P_Y(X)/P_Y(0)}{X}. \label{rl-inverse-polynomial}
\end{align}
The numerator has zero constant term, so $q_Y$ is
a polynomial, of degree at most $|\mathcal S|-1$.
Repeated values of $p_\rho(Y)$ cause no difficulty.
For every $\rho$,
$q_Y(\sigma p_\rho(Y))=(\sigma p_\rho(Y))^{-1}$.
For real $Y$ this polynomial has real coefficients.
An explicit simultaneous choice avoiding the zeros
at $Y-1,Y,Y+1$ is
\[
 Y=2+\max\bigl(\{0\}\cup
          \{|\rho_i|:\rho\in\mathcal S,\ 1\leq i<k\}\bigr).
\]

Let $W_\pi$ be the left coefficient space and set
$\mathcal H=\operatorname{Hom}(W_{\pi'},W_\pi)$ with
the Hilbert--Schmidt inner product.
For $g\in\mathbf G^+$ define the endomorphism
\begin{equation}\label{rl-zero-coefficient}
       M_Y(g):A\longmapsto
          \mathcal G(g)^{-1}A\,q_Y(\sigma C_g(Y)).
\end{equation}

\begin{theorem}\label{rl-regular-inverse}
The coefficient \eqref{rl-zero-coefficient} is smooth
on $\mathbf G^+$ and has the following properties.

\begin{enumerate}
\item On a radial cell it realizes exactly the signed
terminal inverse multiplied by the reciprocal full
Gram polynomial:
\[
 M_Y(b(c)m)\Phi(c,m)
    =\frac1{|a(c)|^2G(m)}
       \left(\kappa\bigl(
            (\sigma p_\rho(Y))^{-1}\bigr)\Phi\right)(c,m).
\]
Thus it is the actual coefficient used for the product
of this inverse and the reciprocal Gram factor.
\item It commutes with the full compact and unipotent
action on coefficient functions.  More explicitly,
for the elements in Lemma~\ref{rl-cofactor-covariance},
\[
 M_Y(g)\bigl(\pi(h_L)A\pi'(h_R)^{-1}\bigr)
    =\pi(h_L)\bigl(M_Y(g')A\bigr)\pi'(h_R)^{-1}.
\]
\item Its pointwise Hilbert--Schmidt adjoint is smooth
and has the same derivative growth bounds as $M_Y$.
For real $Y$ it equals $M_Y$.
\item For every $\varepsilon>0$ and integer $m\geq0$,
there is a finite constant $C_{\varepsilon,m}$ such that
every word $D_1\cdots D_j$ of at most $m$ left or right
Lie derivatives in matrices whose entries have absolute
value at most one satisfies
\begin{equation}\label{rl-growth}
 \|D_1\cdots D_j M_Y(g)\|
       \leq C_{\varepsilon,m}
                      (1+\|g\|_{\rm F})^{(r+1)m}
         \quad(\mathcal G(g)\geq\varepsilon).
\end{equation}
The same bound holds for its adjoint.
The constant depends on the fixed Hermitian data,
the finite representations, $Y,\sigma,\varepsilon,m$,
and not on $g$ or on the selected unit Lie directions.
\end{enumerate}
\end{theorem}

\begin{proof}
The finite decomposition of Theorem~\ref{rd-constructed-kappa}
and \eqref{rl-terminal-match} identify
$q_Y(\sigma C_g(Y))$ on the evaluated radial image
with the specified inverse.
No assertion that it is the inverse on unused
right constituents is needed.  This proves the first
part, using \eqref{rl-cell-Gram}.
Smoothness follows from \eqref{rl-rational-matrix}.
The covariance of $C_g$ and invariance of $\mathcal G$
give the second part by moving $\pi'(h_R)^{-1}$
through the polynomial.
The adjoint of right multiplication by a matrix $B$
on the Hilbert--Schmidt space is right multiplication
by $B^*$.  The adjoint formula in
Lemma~\ref{rl-compressed-capelli} therefore proves
the third part; the final assertion for real $Y$
also uses that $q_Y$ has real coefficients.

Here is a direct derivative estimate, which avoids
differentiating any orthonormal frame.  Write
$d_0=r+1$, $n=2r+k+1$, and use the fixed polynomial
extension of $\Omega$ from Lemma~\ref{rl-cofactor-norm}.
There is a constant $C_\Omega$ such that any word
$L$ of $a$ left or right unit Lie derivatives obeys
\begin{equation}\label{rl-polynomial-derivative}
               \|L\Omega(g)\|
       \leq C_\Omega(nd_0)^a(1+\|g\|_{\rm F})^{d_0}.
\end{equation}
Indeed each such vector field replaces an entry
of $g$ by a sum of at most $n$ entries with coefficients
of absolute value at most one, and likewise for a
conjugate entry.  On a monomial of degree $d_0$
there are $d_0$ choices of a differentiated factor.
Thus each derivative multiplies the coefficient
sum by at most $nd_0$, without increasing the
degree.  The Euclidean sum of the component coefficient
norms gives a possible $C_\Omega$.

For $v\in\mathbb C^k\setminus\{0\}$ set
\[
 \mathcal F_Y(v):A\longmapsto
       \frac{1}{T'\|v\|^2}
          A\,q_Y\bigl(\sigma C_{v/\|v\|}(Y)\bigr).
\]
This is a smooth function with values in
$\operatorname{End}(\mathcal H)$, and
$\mathcal F_Y(tv)=t^{-2}\mathcal F_Y(v)$ for real $t>0$.
Consequently
\[
 \|D^\ell\mathcal F_Y(v)\|
       \leq M_\ell\|v\|^{-2-\ell},\qquad
 M_\ell=\max_{\|u\|=1}\|D^\ell\mathcal F_Y(u)\|<\infty.
\]
The finiteness follows from smoothness and compactness
of the unit sphere.  These formulas define the
constants entirely in the fixed finite coefficient space.
By \eqref{rl-omega-norm},
$M_Y(g)=\mathcal F_Y(\Omega(g))$, and
$\|\Omega(g)\|\geq\delta=\sqrt{\varepsilon/T'}$
on the set under consideration.

The iterated chain rule for a word of $j$ vector fields
expresses its derivative of $\mathcal F_Y(\Omega)$
as a sum over partitions of those $j$ positions.
A partition with $\ell$ blocks contributes
$D^\ell\mathcal F_Y$ applied to $\ell$ derivatives
of $\Omega$, using the original order within each
block.  This statement follows by induction:
the next derivative either differentiates
$\mathcal F_Y$ or one of its arguments.
It does not require the vector fields to commute.
Using \eqref{rl-polynomial-derivative}, its norm is
at most
\[
 (nd_0)^j
 \sum_{\ell=0}^j
    \left\{\begin{matrix}j\\\ell\end{matrix}\right\}
       M_\ell\delta^{-2-\ell}C_\Omega^\ell
                    (1+\|g\|_{\rm F})^{d_0\ell},
\]
where the braces count partitions into $\ell$ nonempty
blocks, with the usual value $1$ at $(j,\ell)=(0,0)$.
Take the maximum of the corresponding fixed constants
over $0\leq j\leq m$, and use $\ell\leq j\leq m$.
This proves \eqref{rl-growth}.
Taking an adjoint is a fixed real-linear isometry
on the finite coefficient space, so the same argument
and bound apply to the adjoint coefficient.
\end{proof}

\subsection{Preservation of the actual rapidly decreasing class}

For $\varepsilon>0$, let $\mathscr S_\varepsilon$
consist of the smooth $\mathcal H$-valued functions
on $\mathbf G$ which vanish on $\mathcal G<\varepsilon$
and for which every mixed left and right Lie derivative
decays faster than every power of $1+\|g\|_{\rm F}$,
uniformly over unit Lie directions.
This is the coefficient class in the source's
smoothness, decay and support hypotheses.
Its formulation with the norm $\|g\|_{\rm F}$ alone
is equivalent: $|\det g|=1$, so Hadamard's determinant
inequality implies $\|g\|_{\rm F}\geq1$.
Hence $1+\|g\|_{\rm F}\leq2\|g\|_{\rm F}$.

\begin{corollary}\label{rl-class-preserved}
Extend $M_YF$ by zero on $\mathcal G\leq0$.
Both $F\mapsto M_YF$ and its pointwise adjoint
preserve $\mathscr S_\varepsilon$.
They are zero-order differential operators with
smooth coefficients on $\mathbf G^+$, they commute
with the source compact and unipotent actions, and
all their coefficient derivatives have polynomial
growth on $\mathcal G\geq\varepsilon$.
\end{corollary}

\begin{proof}
The functions in question vanish on the neighborhood
$\mathcal G<\varepsilon$ of $\mathcal G=0$,
so their zero extensions are smooth there.
On $\mathbf G^+$ use the product rule.
For a derivative of order at most $m$, every
coefficient factor is bounded by
\eqref{rl-growth}.  The differentiated $F$ can be
bounded by an arbitrarily large negative power;
choosing that power larger by $(r+1)m$ gives the
required decay.  The support is not enlarged.
Covariance and the adjoint properties were proved
in Theorem~\ref{rl-regular-inverse}.
\end{proof}

The exponent $(r+1)m$ in \eqref{rl-growth} is sufficient
for the actual qualitative coefficient condition.
It also improves the exponent $3(r+1)m$ used in
the corresponding source provider for this zero-order
factor.  The constants are not claimed to be uniform
as $Y$ approaches a forbidden spectral value.
They can be chosen uniformly on a fixed compact set
of $Y$ avoiding those finitely many values, since
the coefficients of $q_Y$ and all sphere derivatives
above then vary continuously on compact sets.

\begin{example}\label{rl-vanishing-middle-minor}
The domain $\mathbf G^+$ really contains points
where a radial coordinate is unavailable.
Take $k=r=1$, a basis $(e,a,z,e^\vee)$ with
$h(e,e^\vee)=1$, $h(a,a)=1$, $h(z,z)=-1$,
and all other required pairings zero.
Set $z'=z$ and define
\[
 ge=a+z,\quad
 ga=\tfrac12e+e^\vee,\quad
 gz=\tfrac12e-e^\vee,\quad
 ge^\vee=\tfrac12(a-z).
\]
The displayed images have exactly the same pairings
as the basis, so $g\in\mathbf G$.
Here
\[
             Z(g)=\begin{pmatrix}0&-1\\-1&0\end{pmatrix},
       \qquad d(g)=0,\quad D(g)=-1,\quad\mathcal G(g)=1.
\]
Thus a proposed full-group lift involving $d(g)^{-1}$
would be singular at an allowed point.
The cofactor coefficient above is regular there:
$L(g)=e$, $y(g)=a+z$, and $\Omega(g)=1$.
This is not a counterexample to the source, whose
construction uses the regular Gram denominator.
\end{example}

\subsection{What this completes}

The preceding construction proves the actual
zero-order spectral factor, rather than assuming
a smooth matrix field with suitable bounds.
The polynomial cofactor identity proves its full-group
definition; the radial-cell calculation proves that
it is the correct operator; and covariance and
\eqref{rl-growth} prove the analytic coefficient
conditions needed to compose it with other
differential operators.

The remaining full-group operators in the Fourier
identity include two weighted radial transfer
operators and terminal off-diagonal contractions.
Their restriction identities and formal-adjoint
passage through the Jacquet transform are not
proved merely by the zero-order construction here.
Likewise the integrated Jacquet pairing and the
actual local uniform $L^2$ bounds remain separate.
No full weighted Fourier decay theorem is asserted
until those remaining identities are established.


